\documentclass[10pt]{article}

\usepackage[utf8]{inputenc}

\usepackage[T1]{fontenc}

\usepackage{amsmath}

\usepackage{amsfonts}

\usepackage{amssymb}

\usepackage[version=4]{mhchem}

\usepackage{stmaryrd}

\usepackage{graphicx}

\usepackage[export]{adjustbox}

\graphicspath{ {./images/} }

\usepackage{fixltx2e}



\begin{document}

ЭРЛАНГЕНСКАЯ ПРОГРАММА - единая точка зрения на различные геометрии (вапр., евклидову, аффинную, проективную), сформулированная вцервые Ф. Клейном (F. Klein) на лекции, прочитанной в 1872 в ун-те г. Эрланген (Германия) и напечатанной в том же году под назв. «Сравнительное обозрение новейших геометрических исследований».



Сущность Э. п. состоит в следующем. Как известно, евклидова геометрия рассматривает те свойства фигур, к-рые не меняются при движениях; равные фигуры определяются как фигуры, к-рые можно перевести одну в другую движением. Но вместо движений можно выбрать какую-нибудь иную совокупность геометрич. преобразований и объявить "равными» фигуры, получающиеся одна из другой с помощью преобразований этой совокупности; это приводит к иной «геометрии», изучающей свойства фигур, не меняющиеся при рассматриваемых преобразованиях. Введенное «равенство» должно удовлетворять следующим трем естественным условиям: 1) каждая фигура $F$ "равна» сама себе; 2) если фигура $F$ "равна» фигуре $F^{\prime}$, то и $F^{\prime}$ "равна» $F$; 3) если фигура $F$ «равна» $F^{\prime}$, а $F^{\prime}$ «равна» $F^{\prime \prime}$, то и $F$ "равна" $F^{\prime \prime}$. Соответственно әтому надо потребовать, чтобы рассматриваемая совокупность преобразований была группой. Теория, к-рая изучает свойства фигур, сохраняющихся при всех преобразованиях данной групшы, наз. геометрией этой групшы.



Выбор различных групп преобразований приводит к разным геометриям. Так, рассмотрение групшы движений приводит к обычной (евклидовой) геометрии; замена движения аффинными преобразованиями пли проективными преобразованиями приводит к аффинной соответственно проективной геометрии. Основываясь на идеях А. Кэли (A. Cayley), Ф. Клейн показал, что принятие за основу групшы проективных преобразований, переводящих в себя нек-рый круг (или произвольвое кович. сечение), приводит к неевклидовой геометрии Лобачевского. Ф. Клейн ввел в рассмотрение довольно широкий круг других геометрий, определяемых подобным же образом.



Э. п. не охватывает нек-рых важных разделов геометрии, напр. риманову геометрию. Однако Э. п. имела для дальнейшего развития геометрии существенное стимулирующее значение.



Лum.: [1] К л е й н Ф., Сравнительное обозрение новейших геометрических исследований («Әрлангенская программа»), в кн.: Об основаниях геометрии, М., 1956; [2] е г о ж е, Элементарная математика сточки зрения высшей, пер. с нем., 2 изд., т. 2, М.- J., 1934; [3] е г о ж е, Высшая геометрия, пер. с нем., М.- Л., 1939; [4] Е фі и м о в Н. В., Высшая геометрия,



\begin{center}

\includegraphics[max width=\textwidth]{2023_07_09_cfb98be8e8d608371d7cg-1(1)}

\end{center}



ӘРМИТА ИНТЕРПОЛЯЦИОННАЯ ФОРМУЛА форма записи многочлена $H_{m}(x)$ степени $m$, репающего задачу интерполирования функции $f(x)$ и ее производных в точках $x_{0}, x_{1}, \ldots, x_{n}$, т. е. удовлетворяющего условиям:



\[

\begin{aligned}

& H_{m}\left(x_{0}\right)=f\left(x_{0}\right), H_{m}^{\prime}\left(x_{0}\right)=f^{\prime}\left(x_{0}\right), \ldots \\

& \ldots, H_{m}^{\left(\alpha_{0}-1\right)}\left(x_{0}\right)=f^{\left(\alpha_{0}-1\right)}\left(x_{0}\right)

\end{aligned}

\]



\begin{center}

\includegraphics[max width=\textwidth]{2023_07_09_cfb98be8e8d608371d7cg-1}

\end{center}



\[

\begin{aligned}

& \ldots, H_{m}^{\left(\alpha_{n}-1\right)}\left(x_{n}\right)=f^{\left(\alpha_{n}-1\right)}\left(x_{0}\right) \\

& m=\sum_{i=0}^{n} \alpha_{i}-1

\end{aligned}

\]



Э. и. ф. может быть записана в виде:



\[

\begin{gathered}

H_{m}(x)=\sum_{i=0}^{n} \sum_{j=0}^{\alpha_{i}-1} \sum_{k=0}^{\alpha_{i}-1} f^{(j)}\left(x_{i}\right) \times \\

\times \frac{1}{k !} \frac{1}{j !}\left[\frac{\left(x-x_{i}\right)_{i}}{\Omega(x)}\right]_{x=x_{i}}^{(k)} \frac{\Omega(x)}{\left(x-x_{i}\right)^{\alpha_{i} \cdots j-k}},

\end{gathered}

\]



где $\Omega(x)=\left(x-x_{0}\right)^{\alpha_{0}}\left(x-x_{1}\right)^{\alpha_{1}} \ldots\left(x-x_{n}\right)^{\alpha_{n}}$.



Лит.: [1] Б е р е з и н И. С., Ж и д к о в Н. П., Методы вычислений, 3 изд., т. 1, М., 1966 . ЭРМИТА МНОГОЧЛЕНЫ, м н о г о ч л е ны Ч е6 ы ші е в а - Э р м и т а,- многочлены, ортогональные на ивтервале $(-\infty, \infty)$ с весовой функцией $h(x)=$ $=\exp \left(-x^{2}\right)$. Стандартизованные Э. м. определяются Родрига формулой



\[

H_{n}(x)=(-1)^{n} e^{x^{2}}\left(e^{-x^{2}}\right)^{(n)} .

\]



Наиболее употребительны формулы



\[

\begin{gathered}

H_{n+1}(x)=2 x H_{n}(x)-2 n H_{n-1}(x), \\

H_{n}^{\prime}(x)=2 n H_{n-1}(x), \\

H_{n}(x)=\sum_{k=0}^{\left[\frac{n}{2}\right]} \frac{(-1)^{k} n !}{k !(n-2 k) !}(2 x)^{n-2 k}, \\

\exp \left(2 x w-w^{2}\right)=\sum_{n=0}^{\infty} \frac{H_{n}(x)}{n !} w^{n} .

\end{gathered}

\]



Первые Э. м. имеют вид



\[

H_{0}(x)=1, H_{1}(x)=2 x, H_{2}(x)=4 x^{2}-2 \text {, }

\][



H\_\{3\}(x)=8 x\^{}\{3\}-12 x \textbackslash text \{, \}

]



$H_{4}(x)=16 x^{4}-48 x^{2}+12, \quad H_{5}(x)=32 x^{5}-160 x^{3}+120 x, \ldots$



Многочлен $H_{n}(x)$ удовлетворяет диффференциальному уравнению $y^{\prime \prime}-2 x y^{\prime}+2 n y=0$.



Ортонормированные Э. м. определяются равенством



\[

\hat{H}_{n}(x)=\frac{H_{n}(x)}{\sqrt{n ! 2^{n \sqrt{\pi}}}} .

\]



Э. м. с единичным старшим коэффициентом имеют вид



\[

\tilde{H}_{n}(x)=\frac{1}{2^{n}} H_{n}(x)=\frac{(-1)^{n}}{2^{n}} e^{x x^{2}}\left(e^{-x}\right)^{(n)} .

\]



Ряды Фурье по Э. м. внутри интервала $(-\infty, \infty)$ аналогичны тригонометрич. рядам Фурье.



В математич. статистике и теории вероятностей применяются Э. м., соответствующие весовой функции



\[

h(x)=\exp \left(-x^{2} / 2\right) .

\]



Определение Э. м. встречается у П. Лапласа [1]. Подробное исследование этих многочленов опубликовал П. Л. Чебышев в 1859 (см. [2]). Затем эти многочлены изучал II. Эрмит [3]. В. А. Стеклов [4] доказал плотность множества всех многочленов в пространстве функций, квадрат к-рых интегрируем с весом $h(x)=$ $=\exp \left(-x^{2}\right)$ на всей оси.



См. также Классические ортогональные многочлены. Jum.: [1] L a p l a c e P. S., “Mém. classe sci. math., phys. inst. France», 1810 , p. $279-347 ;$ [2] ч е 6 Ыш е в П. Ј., Полн. собр. соч., т. 2, M. - J., 1947, с. $335-41$; [3] H e r m i t e C h. «C. r. Acad. sci.», 1864, t. 58, p. $93-100,266-73$; [4] С т е кло в B. А., "Изв. АН», 1916, т. 10 , с. $403-16 ;$ [5] С у eт и н П. К., Классические ортогональные многочлены, 2 изд.,



ӘРМИТА ПРЕОБРАЗОВАНИЕ - интегральное преобразование вида



\[

f(n)=H\{F(x)\}=

\][



=\textbackslash int\_\{-\textbackslash infty\}\^{}\{\textbackslash infty\} e\textsuperscript{\{-x}\{2\}\} H\_\{n\}(x) F(x) d x, n=0,1,2, ···,

]



где $H_{n}(x)-$ Эрмита многочлен степени $n$. Формула обращения имеет вид



$F(x)=\sum_{n=0}^{\infty} \frac{1}{\sqrt{\pi}} \frac{f(n)}{2^{n} n !} H_{n}(x)=H^{-1}\{f(n)\},-\infty<x<\infty$, если ряд сходится. Э. п. сводит операцию



\[

R[F(x)]=e^{x^{2}} \frac{d}{d x}\left[e^{-x^{2}} \frac{d}{d x} F(x)\right]

\]





\end{document}